% -----------------------------------------------------------------------------
% Theorem styles

% The `nl` style sets the body in italics and breaks the line after the header
\newtheoremstyle{nl}% name
{}% space above
{}% space below
{\itshape}% body font
{}% indent
{\bfseries}% head font
{}% head punctuation
{\newline}% after head space
{\thmname{#1}\thmnumber{ #2}\thmnote{ (#3)}.}% head spec

% The `it` style is the conventional style that sets the body in italics and in
% the same line as the header.
\newtheoremstyle{it}% name
{}% space above
{}% space below
{\itshape}% body font
{}% indent
{\bfseries}% head font
{}% head punctuation
{0.5em}% after head space
{\thmname{#1}\thmnumber{ #2}\thmnote{ (#3)}.}% head spec

% The `rm` style sets the body in roman and in the same line as the header
\newtheoremstyle{rm}% name
{}% space above
{}% space below
{\rmfamily}% body font
{}% indent
{\bfseries}% head font
{}% head punctuation
{0.5em}% after head space
{\thmname{#1}\thmnumber{ #2}\thmnote{ (#3)}.}% head spec

% To customize the theorem styles, set the following options before loading
% this file:
\providecommand{\theoremTheoremStyle}{it}
\providecommand{\exampleTheoremStyle}{rm}

% -----------------------------------------------------------------------------
% Theorem definitions

% Two helper commands
\newcommand{\newTheoremWithOwnCounter}[2]{
    % `#1` is both the environment name and the lowercase reference name
    % `#2` is the capitalized reference name
    \newtheorem{#1}{#2}
    \newtheorem*{#1*}{#2}
    \zcRefTypeSetup{#1}{name-sg=#1,Name-sg=#2,name-pl=#1s,Name-pl=#2s}
}
\newcommand{\newTheoremWithSharedCounter}[3]{%
    % `#1` is both the environment name and the lowercase reference name
    % `#2` is the capitalized reference name
    % `#3` is the counter
    \newtheorem{#1}[#3]{#2}
    \newtheorem*{#1*}{#2}
    \zcRefTypeSetup{#1}{name-sg=#1,Name-sg=#2,name-pl=#1s,Name-pl=#2s}
    \AddToHook{env/#1/begin}{\zcsetup{countertype={#3=#1}}}
}

% The following are typeset in the `it` style
\theoremstyle{\theoremTheoremStyle}
\newTheoremWithOwnCounter{theorem}{Theorem}
\newTheoremWithSharedCounter{proposition}{Proposition}{theorem}
\newTheoremWithSharedCounter{lemma}{Lemma}{theorem}
\newTheoremWithSharedCounter{corollary}{Corollary}{theorem}
\newTheoremWithSharedCounter{conjecture}{Conjecture}{theorem}
\newTheoremWithOwnCounter{assumption}{Assumption}
\newTheoremWithOwnCounter{claim}{Claim}
\newTheoremWithOwnCounter{definition}{Definition}
\newTheoremWithOwnCounter{remark}{Remark}

% The following are typeset in the `rm` style
\theoremstyle{\exampleTheoremStyle}
\newTheoremWithOwnCounter{example}{Example}
\newTheoremWithOwnCounter{exercise}{Exercise}

% The `exampleQed` environment adds a triangle QED symbol to the `example`
\newenvironment{exampleQed}{%
    \renewcommand{\qedsymbol}{\( \triangleleft \)}%
    \pushQED{\qed}\example
}{\popQED\endexample}

% -----------------------------------------------------------------------------
% Proof environments

% Set the proof header in bold
\renewenvironment{proof}[1][\proofname]{\par
    \pushQED{\qed}%
    \normalfont \topsep6\p@\@plus6\p@\relax
    \trivlist
    \item[\hskip\labelsep\bfseries #1\@addpunct{.}]\ignorespaces
}{%
    \popQED\endtrivlist\@endpefalse
}

% A `subproof` environment
\newenvironment{subproof}[1][\proofname]{%
    \renewcommand{\qedsymbol}{\( \triangleleft \)}%
    \begin{proof}[\sc #1]%
}{%
    \end{proof}%
}
